\subsection{Integrated answers to the research questions}
Table~\ref{tab:rq} connects the five phases without averaging incompatible evidence. The historical memory gain motivates a model-design question; it is not added to the destination gain to obtain a larger effect. Similarly, the raw Alibaba download volume documents provenance and processing effort, while the inferential support is thirty-six targets per service. The numerical audit establishes consistency between saved predictions and reported errors, not independence from every unobserved preprocessing decision made by the data publisher.
\begin{table*}[!t]
\caption{Evidence-based answers. Directions are observed within the specified experiment, with the four primary sign tests reported separately.}\label{tab:rq}
\centering\footnotesize
\begin{tabularx}{\linewidth}{p{.08\linewidth}XXX}
\toprule RQ & Main evidence and direction & Scope & Limitation \\\midrule
1 & Original loss curves generally increase; prolonged blocks raise A3 error & Fixed traces and synthetic generators & Real outages and every service are not established \\
2 & Historical E5 improves memory, worsens CPU; fixed DLinear is affine & E/C/A parameterizations separately & Selection and optimization differ \\
3 & A3 does not beat A0 aggregate; E descriptors give mixed results & Independently tuned families & No orthogonal causal descriptor effect \\
4 & A3 uses15 vs26/74 parameters; lower measured forward cost & B1 resident tensors on measured CPU/CUDA & No universal accuracy/energy advantage \\
5 & PatchTST/NLinear lowest primary means; persistence best external MASE & Calibrated destination32 and one Alibaba day & No independent-period or population validation \\
6 & Historical crossed intervals widen; A3/E5 memory10/10 seeds & Fixed services, masks and selected configurations & Service concentration and temporal dependence remain \\
\bottomrule
\end{tabularx}
\end{table*}
